\documentclass[11pt]{article}
\usepackage{amsmath,amssymb,amsthm,hyperref}
\newtheorem{theorem}{Theorem}
\newtheorem{lemma}[theorem]{Lemma}
\title{Rational interval quadrature as a corollary of Hilbert--Kamke theory}
\author{Literature consequence, not a claim of a new upper-bound theorem}
\date{30 September 2026}
\begin{document}
\maketitle

\begin{theorem}\label{main}
There is an absolute constant $C$ such that for every $r\ge2$ and every
integer $K\ge Cr^3 2^r$, rational nodes $x_1,\ldots,x_K\in(0,1)$ exist
with
\[
 \frac1K\sum_{i=1}^K x_i^j=\frac1{j+1}\qquad(0\le j\le r).
\]
This follows from classical Hilbert--Kamke local-solubility results and
the rational-node circle-method formulation of Cui--Xia--Xiang.
\end{theorem}

\paragraph{The classical local input.}
Let $V=(h^j)_{1\le j,h\le r}$ and $N\in\mathbb Z^r$.
Arkhipov's Theorem 3 in \cite{ark85}, printed p.~34, supplies a
positive uniform lower bound for the singular series if $Vt=N$ is
solvable in $t\in\mathbb Z^r$ and
\[
 s>\min\bigl(r^2 2^{2r-1},\ 3r^3 2^r-r\bigr).
\]
In particular the congruences
\[
 \sum_{i=1}^s y_i^j\equiv N_j\pmod{p^a}\qquad(1\le j\le r)
\]
are solvable for every prime $p$ and every $a\ge1$. The singular series
is the product of the nonnegative local densities; absence of a
solution modulo one prime power would make that local density zero.
Compactness of $\mathbb Z_p^s$ therefore gives a simultaneous solution
in $\mathbb Z_p^s$ for each $p$.

\begin{proof}[Proof of Theorem~\ref{main}]
Set $c_j=K/(j+1)$ and $B=|\det V|$. Choose
\[
 P_0=B(r+1)!.
\]
For every positive multiple $P$ of $P_0$, each number
$N_j=c_jP^j$ is an integer divisible by $B$. Thus
$V^{-1}N=(\operatorname{adj}V)N/\det V$ is an integer vector.
Arkhipov's local input, applied to this $N$ with $s=K$, gives a solution
$y\in\mathbb Z_p^K$ for every $p$. Dividing by $P$ gives a solution of
$\sum_i x_i^j=c_j$ in $(P^{-1}\mathbb Z_p)^K$ for every $p$.

The real interval Chebyshev-quadrature theorem of Gilboa--Peled gives a
real equal-weight rule of degree $2r+2$ for every sufficiently large
$K\ge C_0r^2$. Such a rule has at least $r+2$ distinct support points,
by integrating the square of the polynomial vanishing on its support.
There are therefore at least $r$ distinct interior nodes. The first $r$
moment Jacobian on these nodes is invertible. By the implicit function
theorem, move any endpoint nodes inward and correct these $r$ pivots.
This gives a fixed real solution $u\in(0,1)^K$ of the first $r$ moment
conditions, with $r$ distinct interior coordinates.

For all sufficiently large multiples $P$ of $P_0$, this real solution
lies in $[P^{-1},1-P^{-1}]^K$ and has positive separation quantity
$\Delta_u$ in \cite[Theorem 5.2(v)]{cxx}. That theorem's parts (iv) and
(v) give positive uniform lower bounds for its singular series and
singular integral. Part (i) then gives a positive number of rational
solutions in $((0,1)\cap P^{-1}\mathbb Z)^K$ for sufficiently large $P$.
Its required variable threshold, polynomial in $r$ times $2^r$, is
covered by our choice $K\ge Cr^3 2^r$.
\end{proof}

\paragraph{Attribution and sharper classical information.}
The exponential upper bound is thus a direct consequence of earlier
published theorems. The separate file
\texttt{exponential\_rational\_designs.tex} supplies an independent
quantitative local argument with the weaker threshold $Cr^6 4^r$.
The interval-integral valuation argument giving an exponential lower
bound is a separate statement; the unit-indicator polynomials used in
that argument already occur in Hua, Arkhipov, and Mit'kin.

Mit'kin \cite{mitkin} sharpens the classical Hilbert--Kamke variable
threshold to
\[
 H(r)=\sum_{0\le a\le\lfloor\log_2r\rfloor}
        2^a\bigl(2^{\lfloor r/2^a\rfloor}-1\bigr)
     =2^r+O(2^{r/2})\qquad(r\ge12).
\]
The localized interval argument is now supplied in
\texttt{mitkin\_sharp\_rational\_upper.tex}: for every sufficiently
large $r$, every $K\ge H(r)$ admits a rational interval rule. In
particular $R_r\le2^r+O(2^{r/2})$. The simpler argument above remains
a direct application of Cui--Xia--Xiang with a larger threshold.

\begin{thebibliography}{9}
\bibitem{ark85} G.~I.~Arkhipov,
\emph{On the Hilbert--Kamke problem},
Mathematics of the USSR-Izvestiya 24 (1985), no.~1, 1--47,
\href{https://doi.org/10.1070/IM1985v024n01ABEH001210}{doi:10.1070/IM1985v024n01ABEH001210}.
\bibitem{cxx} Zhen Cui, Jiacheng Xia, and Ziqing Xiang,
\emph{Rational designs}, Advances in Mathematics 352 (2019), 541--571,
\href{https://doi.org/10.1016/j.aim.2019.06.012}{doi:10.1016/j.aim.2019.06.012}.
\bibitem{gp} Shoni Gilboa and Ron Peled,
\emph{Chebyshev-type quadratures for doubling weights},
Constructive Approximation 45 (2017), 193--216,
\href{https://arxiv.org/abs/1507.01505}{arXiv:1507.01505}.
\bibitem{mitkin} D.~A.~Mit'kin,
\emph{An estimate for the number of terms in the Hilbert--Kamke problem},
Mathematics of the USSR-Sbornik 57 (1987), no.~2, 561--590,
\href{https://doi.org/10.1070/SM1987v057n02ABEH003087}{doi:10.1070/SM1987v057n02ABEH003087}.
\end{thebibliography}
\end{document}
